% ===========================================================================
\chapter{Negation, in layers}\label{ch:negation}

Many rules are about what is \emph{not} so. Look at \code{access.dl}:

\begin{shell}
can_read(U, F) :- user(U), file(F), not banned(U), not secret(F).
can_read(U, F) :- manager(U), file(F), secret(F).

refused(U, F) :- user(U), file(F), not can_read(U, F).
\end{shell}

Anyone not banned may read a file that is not secret; managers may read
secret files too; everything else is refused. The join has handled
\code{not} since Chapter~\ref{ch:joins}: \code{not can_read(U, F)} holds when
that fact is not in the database. But \emph{when} is it not in the
database?

\section{When negation goes wrong}

With positive rules, the order in which facts are derived does not matter:
a fact derived late is as good as one derived early, and nothing is ever
taken back. A negated atom breaks that. \code{not can_read(ann, payroll)}
is true at the start, before any \code{can_read} facts exist, and becomes
false once the second rule has run. If \code{refused} is derived in
between, it is derived from something that later stops being true, and
facts are never removed.

Here is the access program evaluated naively, with the \code{refused} rule
written first, so that it runs before \code{can_read} has any facts:

\begin{tryit}{try_wrong.nex}
intern Parser
intern Evaluate

match parse_program("user(ann). user(bo). file(report). banned(bo).
refused(U, F) :- user(U), file(F), not can_read(U, F).
can_read(U, F) :- user(U), file(F), not banned(U).") of
  Ok as ok then
    do
      let db := create Database.make
      across ok.value.facts as f do
        db.add(f)
      end
      (create Naive_Evaluator.make).run(ok.value.rules, db)
      across ["refused", "can_read"] as p do
        across db.facts_of(p) as f do
          console.print_line(f.text)
        end
      end
    end
  Err as e then console.print_line(e.error)
end
\end{tryit}
\tryoutput{try_wrong}

Ann is refused the report, and also allowed to read it. The first round
derived \code{refused(ann, report)} while \code{can_read} was still empty,
and nothing ever withdrew it.

\section{A precondition}

The cure is to finish a predicate before anything negates it: evaluate all
of \code{can_read} first, then the rules that say \code{not can_read}. A
negated predicate must be \emph{settled}, complete, before the rules that
negate it run. That is a property of the rules handed to one call of
\code{run}: none of them may negate a predicate that another of them
defines.

The rule for suggestions in \code{social.dl} already has this property.
\code{suggest} negates only \code{follows}, and no rule defines
\code{follows}, so it is settled from the start. That is why the suggestions
in the earlier chapters were right before any of this existed.

\begin{append}{terms.nex}
-- The predicate a negated literal denies, or "" for any other literal.
function negated_predicate(l: Literal): String
do
  match l of
    Negated(atom) then result := atom.predicate
    _             then result := ""
  end
end

-- Does no rule negate a predicate that one of `rules` defines? Then each
-- negated predicate is already complete when the rules are applied.
function negations_settled(rules: Array[Rule]): Boolean
do
  result := true
  across rules as r do
    across r.body as l do
      let p := negated_predicate(l)
      across rules as other do
        if p /= "" and other.head.predicate = p then
          result := false
        end
      end
    end
  end
end
\end{append}

The property becomes a precondition of every evaluator:

\begin{tightaddbefore}{evaluate.nex}{to run, before deferred}{  deferred|  ensure|    fixpoint: is_fixpoint(rules, db)}
    negations_settled: negations_settled(rules)
\end{tightaddbefore}

\section{Strata}

Now a program has to be evaluated in layers, called \emph{strata}. Give each
predicate a level. A rule's head must be at least as high as every
predicate the rule uses positively, and higher than every predicate it
negates. Then evaluate the rules level by level, lowest first: by the time
a level is evaluated, every predicate it negates is in a lower level, and
complete.

The levels can be found by starting everything at 0 and raising heads until
every rule is satisfied. That fails for one kind of program, one that defines
a predicate through its own negation:

\begin{shell}
p(X) :- q(X), not p(X).
\end{shell}

Here \code{p} must be higher than itself, and the levels rise forever. No
evaluation order makes sense of such a program, so the engine refuses it.
Levels can never legitimately rise above the number of predicates, so when
one does, the stratifier stops and reports the rule.

\begin{newfile}{analysis.nex}
-- analysis.nex: checking a whole program before running it, and running
-- it one stratum at a time.

intern Evaluate

-- The stratum of every predicate the rules define: the predicates a rule
-- negates must come in an earlier stratum than its head, and those it
-- uses positively in the same stratum or an earlier one.
class Stratifier
create
  make(rules: Array[Rule])
  do
    level := {}
    problem := ""
    across rules as r do
      level.put(r.head.predicate, 0)
    end
    -- Raise levels until every rule is satisfied. A program whose levels
    -- keep rising past the number of predicates has a cycle through a
    -- negation, and cannot be stratified.
    let changed := true
    from until not changed or problem /= "" do
      changed := false
      across rules as r do
        across r.body as l do
          let needed := required_level(l)
          if needed > level.get(r.head.predicate) then
            level.put(r.head.predicate, needed)
            changed := true
            if needed > level.size then
              problem := r.head.predicate + " depends on its own negation, through the rule: " + r.text
            end
          end
        end
      end
    end
  end
feature
  level: Map[String, Integer]
  -- Why the rules cannot be stratified, or "".
  problem: String

  -- The lowest level a rule's head can have, given one body literal.
  required_level(l: Literal): Integer
  do
    result := 0
    match l of
      Positive(atom) then result := level.try_get(atom.predicate, 0)
      Negated(atom)  then result := level.try_get(atom.predicate, 0) + 1
      Compare        then result := 0
    end
  end

  -- The rules, grouped by the level of their heads, lowest first.
  strata(rules: Array[Rule]): Array[Array[Rule]]
  require
    stratified: problem = ""
  do
    result := []
    let top := 0
    across level.values as n do
      top := top.max(n)
    end
    from let k := 0 until k > top do
      let group: Array[Rule] := []
      across rules as r do
        if level.get(r.head.predicate) = k then
          group.add(r)
        end
      end
      if not group.is_empty then
        result.add(group)
      end
      k := k + 1
    end
  ensure
    each_settled: all_settled(result)
  end

  all_settled(groups: Array[Array[Rule]]): Boolean
  do
    result := true
    across groups as g do
      if not negations_settled(g) then
        result := false
      end
    end
  end
end
\end{newfile}

\code{strata}'s postcondition is the whole point: every group it returns
satisfies \code{run}'s precondition. \code{program_problem} gathers every
check a program must pass before it can run, safety and stratification, into
one message, and \code{evaluate_program} loads a program's facts and
evaluates its strata in order. Its own postcondition is the fixpoint of the
whole program.

\begin{continued}{analysis.nex}
-- What is wrong with `p` before it can run, or "" if nothing is.
function program_problem(p: Datalog_Program): String
do
  result := ""
  across p.rules as r do
    if result = "" and not r.is_safe then
      result := r.safety_problem + ", in the rule " + r.text
    end
  end
  if result = "" then
    result := (create Stratifier.make(p.rules)).problem
  end
end

-- Load `p`'s facts into `db` and derive everything its rules imply, one
-- stratum at a time.
function evaluate_program(p: Datalog_Program, db: Database, e: Evaluator)
require
  runnable: program_problem(p) = ""
do
  across p.facts as f do
    db.add(f)
  end
  let stratifier := create Stratifier.make(p.rules)
  across stratifier.strata(p.rules) as group do
    e.run(group, db)
  end
ensure
  fixpoint: is_fixpoint(p.rules, db)
end
\end{continued}

\section{Trying it out}

\begin{tryit}{try_negation.nex}
intern Parser
intern Analysis
intern Query
intern io/Path

function run_file(name: String)
do
  match parse_program(create Path.make(program.working_directory + "/" + name).read_text) of
    Ok as ok then
      do
        let problem := program_problem(ok.value)
        if problem /= "" then
          console.print_line(name + ": " + problem)
        else
          let db := create Database.make
          evaluate_program(ok.value, db, create Semi_Naive_Evaluator.make)
          across ok.value.queries as q do
            console.print_line("?- " + q.text + ".   " + answer_lines(q, db).to_string)
          end
        end
      end
    Err as e then console.print_line(e.error)
  end
end

run_file("access.dl")
match parse_program("q(a). p(X) :- q(X), not p(X).") of
  Ok as ok then console.print_line(program_problem(ok.value))
  Err as e then console.print_line(e.error)
end
\end{tryit}
\tryoutput{try_negation}

Bo is banned, so he is refused both files; Cy is not a manager, so he is
refused the payroll; Ann may read both. And the program that negates itself
is refused, with the rule that does it.
